\section{결과 보고와 판정 원칙}

실증 결과는 재현 스크립트가 생성한 기계판독 가능 표를 원천으로 삼는다. 본문에
직접 입력한 수치와 수작업으로 다듬은 그림은 사용하지 않는다. 결과는 주분석과
행동모형의 두 묶음으로 분리하며, 행동모형의 성패가 주분석 결과의 보고 여부를
결정하지 않는다.

\subsection{주분석의 보고 순서}

첫 번째 표는 승식별 표본 경주 수, 완전 조합 비율, 배당상한과 반올림 빈도를
보고한다. 이어서 3,321개 clean 경주와 15,963개 상한 경주의 출전두수, 연도,
경마장, 상한 조합 수와 상한 직전 비검열 배당의 꼬리분포를 비교한다. 주결과표는
두 패널을 고정한다. Panel A는 clean 경주의 경주균등가중
TV·Jensen--Shannon·log-ratio 거리, calibration과 기준모형 대비 손실 차이의
점추정과 신뢰구간을 제시한다. Panel B는 전체 19,284개 경주에서 TV·MAE의
정확한 하한과 인증된 보수적 상한, 기준모형의 대응 경계를 제시한다. 비선형
JS·log-ratio·calibration의 Panel B 결과는 설명적 민감도 분석으로 구분한다.
두 패널에는 같은 주변화, 가중 방식과 Harville·순열·타경주·균등 기준을 적용한다.
패널의 1차 집계치는 경주균등가중 중앙값이다. 기준모형 우위는 Panel A의 짝지은
TV 개선폭 중앙값과 Panel B의 보수적 개선폭 하한 중앙값을 사용하고, 각각의
경주 단위 부트스트랩 95\% 신뢰구간 하한이 0보다 큰지 보고한다. 절대오차에는
$\tau_{\rm TV}=0.05$를 사전 기준으로 사용하고 0.025와 0.10의 민감도를 제시한다.

해석은 세 단계를 따른다. 먼저 가격 재구성의 절대 오차가 경제적으로 작은지
판단한다. 그 다음 그 오차가 기준모형보다 작은지 확인한다. 마지막으로 연도와
경마장을 바꾸어도 차이가 유지되는지 검토한다. 높은 전체표본 $R^2$만으로
``삼쌍승식이 모든 정보를 담는다''거나 ``공통 상태가격이 존재한다''고 결론
내리지 않는다. 별도 풀의 유동성, 공제, 절사와 참가자 구성은 식
\eqref{eq:cross-market-restriction}에서 벗어나는 경제적 원인일 수 있다.
Panel A의 개선폭 중앙값만 유의하게 양수이고 Panel B의 보수적 개선폭 하한
중앙값은 그렇지 않으면 결론은 clean 표본으로 제한한다. 두 패널이 같은 방향을
지지할 때만 전체 기간·경주로 일반화한 강한 결론을 허용한다.

\subsection{행동모형의 보고 순서}

행동모형 표는 먼저 조건부 순위확률의 표본 밖 로그점수와 calibration을
보고한다. 이 단계가 통과된 뒤에만 M-U, M-R, M-S2와 M-S3의 가격예측 오차를
제시한다. 각 모형은 단승→쌍승, 단승→삼쌍승, 순서형→무순서형,
훈련연도→검증연도의 이동성으로 나누어 평가한다. 표본 안 적합도보다 연도 밖과
leave-one-pool-out 손실을 우선한다.

M-S2 또는 M-S3의 우위는 복합복권 비축약과 양립하는 증거다. 그러나 집계자료로
개인의 위험태도, 신념과 사건 프레임을 동시에 식별했다고 표현하지 않는다.
위험선호 M-U와 축약형 확률가중 M-R의 차이가 함수형 제약에 민감하면 모형순위를
강한 심리적 결론으로 바꾸지 않는다. 풀별 재추정이 필요한 경우에는 참가자
이질성 또는 풀별 가격형성 차이를 우선적인 설명으로 남긴다.

주요 결론은 분석 코드, 표본 포함 규칙, 가중 방식, 확률모형, 기준모형과
검증 fold가 모두 고정된 뒤에만 서술한다. 탐색 과정에서 발견한 패턴은 탐색
결과로 표시하고, 확인 표본이나 시간 분할 검증을 거치지 않은 패턴에는 확정적
언어를 쓰지 않는다. 내부 가격 정합성, 실현 착순 예측과 수익성·후생은 서로
다른 표와 결론으로 분리한다.
